\chapter{Arquitectura y diseño del sistema inteligente}
\label{cap:6}

Este capítulo describe la arquitectura de software y el diseño detallado del sistema inteligente de priorización. En primer lugar, se enuncian los principios de diseño que rigen todo el desarrollo, tales como la explicabilidad y el control humano de la decisión. A continuación, se detalla la organización modular del sistema estructurada en tres capas diferenciadas: extracción textual, modelo predictivo e interpretación normativa. Por último, se describe la especificación de las interfaces y los contratos de datos que garantizan la consistencia del flujo de información a lo largo de toda la canalización.

\section{Finalidad del capítulo}

Este capítulo presenta el diseño técnico y analítico del sistema propuesto. Tras haber fijado el problema, el estado del arte y el marco normativo, aquí se describe cómo se transforma un incidente 112 en una recomendación de prioridad P1--P4 explicable, auditable y revisable por un operador humano. El modelado en tres capas independientes constituye la respuesta técnica a la necesidad de aislar la toma de decisiones estadísticas de la generación lingüística explicativa.

La versión actual del TFM abandona el diseño inicial basado únicamente en reglas duras y \textit{scoring}. Esta decisión responde a la necesidad de dotar al prototipo de una mayor capacidad de generalización sobre registros reales. En consecuencia, la arquitectura final se organiza en tres capas de inteligencia artificial, coordinadas mediante contratos de datos versionados:

\begin{itemize}
    \item \textbf{Capa 1: NLP operativo}. Extrae variables V01--V15 y señales textuales desde el título, descripción, categoría preliminar y contexto básico del incidente.
    \item \textbf{Capa 2: priorización interpretable}. Produce una recomendación P1--P4 mediante RuleFit-lite, entrenado sobre etiquetas académicas generadas por supervisión débil, y comparado con una línea base experta y una ejecución diagnóstica con \texttt{imodels}.
    \item \textbf{Capa 3: explicación normativa}. Genera una explicación breve para el operador mediante LLM local, recuperación aumentada sobre corpus normativo de Castilla y León y herramientas MCP.
\end{itemize}

El resultado no es un sistema autónomo de despacho, sino un sistema de apoyo a la decisión. La prioridad recomendada no sustituye el criterio profesional; se ofrece como una propuesta estructurada, con reglas activadas, probabilidad, confianza, citas legales y registro auditable. Al delegar la responsabilidad final en el operador, se asegura el alineamiento con las directrices regulatorias europeas sobre el uso ético de algoritmos de decisión.

\section{Principios de diseño}

El diseño se apoya en seis principios que conectan los capítulos previos con la implementación.

\textbf{Primero, supervisión humana.} El operador acepta, modifica o rechaza la recomendación. El sistema nunca moviliza recursos por sí mismo ni declara situaciones operativas oficiales.

\textbf{Segundo, explicabilidad desde el núcleo.} La prioridad se calcula mediante un modelo interpretable, no mediante un LLM generativo. El LLM explica una decisión ya producida por Capa 2; no decide la prioridad.

\textbf{Tercero, contratos antes que integración.} Cada capa intercambia objetos Pydantic versionados. Esto reduce ambigüedad, facilita pruebas y permite que los tres bloques de trabajo avancen en paralelo.

\textbf{Cuarto, soberanía del dato.} La inferencia se plantea en local. No se envían incidentes a APIs externas en producción, y la Capa 3 usa un LLM cuantizado ejecutado en hardware propio.

\textbf{Quinto, trazabilidad normativa.} Las variables, reglas y explicaciones se vinculan con normas del corpus CyL: Ley 17/2015, RD 524/2023, Ley 4/2007 CyL, PLANCAL, INFOCAL, INUNCYL, MPCyL, RGPD, LOPDGDD y Reglamento UE 2024/1689.

\textbf{Sexto, prevención de \textit{leakage}.} El sistema no utiliza columnas posteriores a la decisión, como medios movilizados, pacientes atendidos o cierre del incidente. Estas variables pueden analizarse solo como información ex post de evaluación, nunca como entrada de Capa 1 o Capa 2.

\section{Vista general de la arquitectura}

La arquitectura se representa como una cadena de transformaciones controladas. La entrada inicial es un \texttt{IncidentInput}; la Capa 1 produce \texttt{IncidentFeatures}; la Capa 2 produce \texttt{PriorityRecommendation}; la Capa 3 produce \texttt{OperatorRecommendation}; y la decisión final del operador se registra como \texttt{OperatorDecision}. Todo el proceso queda recogido en un \texttt{InferenceLog}. Esta secuencia evita dependencias circulares entre las fases de inferencia. De forma opcional, la interfaz puede proponer la entrada a partir de audio de prueba mediante reconocimiento automático del habla (ASR) con \texttt{faster-whisper}; el resultado debe revisarse y no se presenta como transcripción validada de llamadas reales.

\begin{figure}[htbp]
\centering
\small
\fbox{%
\begin{tabular}{c}
\texttt{Streamlit UI / cliente API} \\
\(\downarrow\) \\
\texttt{FastAPI + contratos Pydantic + orquestador de pipeline} \\
\(\downarrow\) \\
\texttt{Capa 1: IncidentInput -> IncidentFeatures} \\
\(\downarrow\) \\
\texttt{Capa 2: RuleFit-lite + Safety Gate -> PriorityRecommendation} \\
\(\downarrow\) \\
\texttt{Capa 3: RAG ChromaDB + servidor MCP localhost + Ollama local} \\
\(\downarrow\) \\
\texttt{OperatorRecommendation -> SQLite / InferenceLog / OperatorDecision}
\end{tabular}
}
\caption{Arquitectura integrada de tres capas con contratos versionados y componentes locales.}
\label{fig:flujo-logico-sistema}
\end{figure}

Esta separación evita mezclar responsabilidades. La Capa 1 no decide prioridad; estructura el incidente. La Capa 2 no redacta explicaciones normativas extensas; recomienda prioridad y devuelve reglas activadas. La Capa 3 no recalcula el resultado; explica lo recibido, cita normas y advierte de incertidumbres.

\section{Contratos de datos}

Los contratos son la pieza que convierte la arquitectura en un sistema verificable. Cada contrato define campos, restricciones e invariantes. La Tabla~\ref{tab:contratos-cap6} resume los objetos principales.

\begin{table}[htbp]
\centering
\small
\begin{tabular}{|p{3.4cm}|p{4.1cm}|p{6.2cm}|}
\hline
\textbf{Contrato} & \textbf{Momento de uso} & \textbf{Función principal} \\ \hline
\texttt{IncidentInput} & Entrada inicial & Representa texto, categoría, localización, fecha, provincia y operador antes de cualquier inferencia \\ \hline
\texttt{IncidentFeatures} & Salida Capa 1 & Contiene variables V01--V15, señales textuales, confianza y metadatos de extracción \\ \hline
\begin{tabular}[c]{@{}l@{}}\texttt{Priority}\\\texttt{Recommendation}\end{tabular} & Salida Capa 2 & Contiene prioridad P1--P4, probabilidades, reglas activadas, nivel de confianza y modelo usado \\ \hline
\begin{tabular}[c]{@{}l@{}}\texttt{Operator}\\\texttt{Recommendation}\end{tabular} & Salida Capa 3 & Contiene explicación textual, citas legales, recomendaciones no vinculantes y resumen de reglas \\ \hline
\texttt{OperatorDecision} & Retroalimentación humana & Registra prioridad final asignada por el operador y posible divergencia \\ \hline
\texttt{InferenceLog} & Auditoría & Conserva snapshots, latencias, versiones de modelo, herramientas invocadas y hash de entrada \\ \hline
\end{tabular}
\caption{Contratos principales del flujo de inferencia.}
\label{tab:contratos-cap6}
\end{table}

El uso de contratos tiene una consecuencia metodológica importante: la memoria puede justificar el sistema sin depender de detalles accidentales de implementación. Lo relevante no es solo que exista código, sino que cada capa tenga entradas, salidas e invariantes comprobables. Asimismo, esta formalización facilita la validación cruzada y la auditoría interna de los componentes del software.

\section{Entrada del sistema: IncidentInput}

La entrada representa lo que el operador o un cliente externo proporciona en los primeros minutos del incidente. Incluye título, descripción, categoría preliminar, localización opcional, provincia, fecha y operador. En el caso del dataset de Castilla y León, estos campos se derivan de columnas como \texttt{Titulo}, \texttt{DescripcionBlob}, \texttt{FechaIncidente}, \texttt{Localidad} y coordenadas.

El contrato exige dos invariantes básicos. Primero, si existe latitud debe existir longitud, y viceversa. Segundo, el título o la descripción deben contener información textual mínima. Esta restricción evita ejecutar una cadena de inferencia sobre entradas vacías o inútiles.

La entrada también es el primer punto de defensa contra \textit{leakage}. Si el cliente incluye campos como medios movilizados o pacientes atendidos, el backend debe rechazarlos con error explícito. Esta decisión protege la validez del sistema: un motor de primera decisión no puede aprender de información que solo aparece después de decidir.

\section{Capa 1: extracción NLP de variables operativas}

La Capa 1 transforma texto libre y metadatos básicos en una representación estructurada del incidente. Su salida, \texttt{IncidentFeatures}, contiene las variables operativas V01--V15, las señales textuales y metadatos de inferencia.

\subsection{Variables activas y diferidas}

El alcance evaluado activa V01--V07 y V12--V15. Las variables V08--V11 permanecen en el contrato como extensión futura, pero no se entrenan ni se exigen para la evaluación principal. Esta decisión mantiene el sistema reproducible y evita depender de integraciones geoespaciales o meteorológicas todavía no consolidadas.

\begin{table}[htbp]
\centering
\scriptsize
\begin{tabular}{|p{1.1cm}|p{4.4cm}|p{4.1cm}|p{3.5cm}|}
\hline
\textbf{Var.} & \textbf{Nombre operativo} & \textbf{Función en el sistema} & \textbf{Estado evaluado} \\ \hline
V01 & Riesgo vital & Detectar amenaza directa a la vida & Activa \\ \hline
V02 & Número de víctimas estimado & Aproximar magnitud humana & Activa \\ \hline
V03 & Gravedad de lesiones & Distinguir leve, grave o crítica & Activa \\ \hline
V04 & Tipo de incidente normalizado & Conectar texto con taxonomía & Activa \\ \hline
V05 & Población vulnerable & Elevar sensibilidad por colectivos vulnerables & Activa \\ \hline
V06 & Número de llamadas & Captar confirmación o simultaneidad & Activa \\ \hline
V07 & Emplazamiento crítico & Identificar hospitales, colegios, residencias u otros puntos sensibles & Activa \\ \hline
V08--V11 & AEMET, inundabilidad, Seveso & Enriquecimiento contextual avanzado & Diferidas \\ \hline
V12 & Riesgo de propagación & Captar incendios u otros incidentes expansivos & Activa \\ \hline
V13 & Multirriesgo & Detectar concurrencia de amenazas & Activa \\ \hline
V14 & Avisos simultáneos en zona & Aproximar presión operativa y confirmación externa & Activa \\ \hline
V15 & Accesibilidad de recursos & Modular complejidad de intervención & Activa \\ \hline
\end{tabular}
\caption{Variables operativas del sistema y alcance evaluado.}
\label{tab:variables-cap6}
\end{table}

\subsection{Señales textuales}

Además de variables estructuradas, la Capa 1 extrae señales textuales booleanas con confianza asociada. Las señales críticas representan lesiones graves, atrapamiento, intoxicación, fallecimiento o riesgo vital. Otras aportan contexto sobre incendios, accidentes de tráfico, rescates, avisos simultáneos y fenómenos meteorológicos o inundaciones. La correspondencia con los identificadores de software se documenta en el Anexo~\ref{cap:anexof}.

La Capa 1 se implementa mediante NLP simbólico y auditable. La detección semántica
basada en reglas y diccionarios operativos, los patrones de negación, la normalización
y la estimación de confianza permiten reconocer señales de alto impacto, como
atrapamiento, herido grave, incendio, intoxicación o riesgo vital textual. Esta
decisión reduce dependencias externas,
facilita los tests de contrato y evita presentar como validado un modelo neuronal que
no dispone de checkpoint congelado. El uso de modelos transformer para extracción
contextual queda reservado como línea futura, condicionada a entrenamiento, versionado
y evaluación específica.

\section{Capa 2: priorización interpretable con RuleFit}

La Capa 2 consume \texttt{IncidentFeatures} y produce \texttt{PriorityRecommendation}. Es el núcleo decisorio del sistema. Su responsabilidad es asignar una prioridad P1--P4, generar probabilidades por clase, devolver reglas activadas y comunicar nivel de confianza.

\subsection{Escala P1--P4}

La escala P1--P4 es académica y funcional, no oficial. P1 representa incidentes críticos con riesgo vital, lesiones graves, atrapamiento, intoxicación, fallecimiento, propagación grave o amenaza inmediata. P2 representa alta prioridad, potencial de deterioro o afectación significativa. P3 recoge incidentes que requieren gestión o seguimiento sin criticidad inmediata. P4 se reserva para baja urgencia relativa.

La escala no equivale a situaciones operativas 0--3. Una situación operativa es una categoría institucional de mando y coordinación; P1--P4 ordena incidentes individuales dentro del prototipo. Esta distinción evita confundir una herramienta académica de priorización con una declaración jurídica de emergencia.

\subsection{Supervisión débil}

Como el Registro 112 CyL no contiene prioridad oficial P1--P4, el entrenamiento de Capa 2 requiere construir una etiqueta académica. La estrategia es la supervisión débil: cuatro funciones de etiquetado heurísticas y correlacionadas generan votos que se agregan en una etiqueta final con confianza. Todas comparten una puntuación programática de gravedad, con variaciones en el uso de la categoría y en patrones de intensificación, vulnerabilidad, emplazamiento y negación. En la ejecución evaluada no intervienen un modelo NER entrenado, un algoritmo de agrupamiento semántico ni un LLM anotador.

La etiqueta débil no se presenta como verdad institucional. Es un recurso metodológico que permite entrenar y comparar modelos sobre datos reales abiertos. Por ello debe acompañarse de acuerdo entre fuentes de etiquetado, análisis de discrepancias y una ablación anti-circularidad que compare el resultado sin la fuente de reglas heurísticas.

\subsection{RuleFit como modelo principal}

RuleFit se selecciona porque combina aprendizaje estadístico e interpretabilidad. El modelo genera o utiliza reglas derivadas de árboles y las pondera mediante regularización, produciendo un conjunto de reglas con coeficientes inspeccionables \cite{Friedman2008RuleFit}. En este TFM se impone una restricción de parsimonia: el modelo seleccionado debe mantener como máximo 30 reglas activas con peso distinto de cero.

Esta restricción es académica, no estética. Un modelo con cientos de reglas sería difícil de defender ante un operador o un tribunal. En cambio, un conjunto reducido de reglas permite revisar por qué P1 aumenta, qué variables influyen y qué anclajes normativos sostienen la recomendación.

\subsection{Línea base experta, RuleFit-lite e \texttt{imodels}}

La Capa 2 no se evalúa en aislamiento. Se compara con dos referencias. La línea base experta es un conjunto de reglas transparentes basadas en señales y variables V01--V15, con anclaje normativo. Sirve como referencia mínima: RuleFit debe aportar mejora cuantitativa sin perder explicabilidad.

La segunda referencia es \texttt{imodels.RuleFitClassifier}, utilizada como implementación canónica de RuleFit. Dado que la librería trabaja en clasificación binaria, se adopta un esquema one-vs-rest P1--P4. La versión evaluada del TFM selecciona RuleFit-lite porque mantiene la filosofía RuleFit, respeta el límite de reglas activas y obtiene el mejor rendimiento de validación en la configuración ensayada.

La alternativa externa se conserva como contraste diagnóstico de la familia RuleFit. El modelo RuleFit-lite se ajusta sobre un submuestreo estratificado y reproducible de 3.000 registros de los 6.512 disponibles, mientras que \texttt{imodels} se limita a 50 registros por su coste local. La diferencia de tamaños impide atribuir a RuleFit-lite una superioridad computacional en condiciones equivalentes; la selección se apoya en el rendimiento de validación, la interpretabilidad y la reproducibilidad de la configuración efectivamente evaluada.

\subsection{Probabilidad, confianza e invariantes}

La salida de Capa 2 incluye una distribución de probabilidad sobre P1--P4. El contrato exige que las probabilidades sumen 1, que la prioridad recomendada coincida con el argmax y que los incidentes P1/P2 incluyan al menos una regla activada cuando el modelo usado sea RuleFit o la línea base. Este requerimiento de consistencia matemática y lógica impide la generación de clasificaciones contradictorias. Asimismo, obliga a que toda asignación crítica cuente con un sustento explícito y auditable en forma de predicados lógicos.

El nivel de confianza se deriva estrictamente de la probabilidad máxima del vector contractual: alta si \(p_{max} \geq 0{,}80\), media entre 0,60 y 0,80, baja entre 0,40 y 0,60, y desconocida por debajo de 0,40. Esta confianza no mide la gravedad física del suceso. Un incidente puede ser P1 con confianza baja o media si hay indicios críticos pero información fragmentaria; del mismo modo, un caso P4 puede tener confianza alta si las señales disponibles son claras.

\subsection{Puerta de Seguridad}

Para reforzar la alineación conceptual con PLANCAL y la Ley 17/2015 de Protección Civil, el diseño de la Capa 2 incorpora un posprocesador denominado Puerta de Seguridad (\textit{Safety Gate}). Este componente aplica una salvaguarda determinista y trazable sobre la salida del modelo. Kollipara et al. \citeA{Kollipara2025TrustworthyML} describen el uso de puertas de decisión verificables como mecanismo de control frente a errores graves en sistemas de aprendizaje automático de dominios críticos. En este TFM su aplicación constituye una decisión de diseño conservadora, no una garantía certificada de seguridad operacional.

Técnicamente, la Puerta de Seguridad actúa como un filtro duro sobre el argmax probabilístico generado por RuleFit-lite. Su función no es sustituir el modelo interpretable, sino reducir el riesgo de que determinados patrones predecisionales de alta gravedad queden infrapriorizados por ruido textual, escasez de muestras o ambigüedad. Para ello evalúa condiciones deterministas basadas en los criterios documentados del proyecto. Su efecto se comprueba mediante escenarios funcionales, mientras que las métricas comparativas de RuleFit-lite se calculan sin este posprocesamiento para no mezclar el rendimiento del estimador con la salvaguarda.

Conceptualmente, RuleFit-lite produce la prioridad
\[
\hat{y}=\arg\max_{c \in \{P1,P2,P3,P4\}} p(c).
\]
La Puerta de Seguridad evalúa después un conjunto de predicados deterministas \(S_i(x)\) sobre las variables de entrada. Cuando se activa un predicado crítico, la prioridad se fuerza a \(P1\) o \(P2\), se incorpora una regla de seguridad con peso alto y se adjuntan anclajes normativos para la explicación posterior. El vector de probabilidad se redistribuye de forma determinista:

\[
p'(y_{seguro}) = 0{,}97,\qquad p'(c \neq y_{seguro}) = 0{,}01.
\]

La Puerta de Seguridad fuerza P1 ante riesgo vital, fallecido, riesgo vital textual, multirriesgo en emplazamiento crítico, incendio forestal con propagación, incendio con propagación en emplazamiento crítico, intoxicación grave, rescate complejo con baja accesibilidad o múltiples víctimas estimadas. Fuerza P2 cuando existen señales operativas relevantes como herido grave, atrapamiento, intoxicación, rescate activo o episodio meteorológico/inundación en emplazamiento crítico, siempre que no se cumpla un criterio P1. Esta lógica mantiene el principio de prudencia: ante indicadores predecisionales críticos, el sistema prefiere elevar la prioridad y exigir revisión humana.

La regla inyectada por la Puerta de Seguridad se incorpora a \texttt{activated\_rules} con identificadores como \texttt{SAFE-GATE-P1-RIESGO-VITAL} o \texttt{SAFE-GATE-P2-OPERATIVO}. De este modo, Capa 3 puede explicar que la recomendación no procede solo del argmax estadístico, sino también de una salvaguarda determinista trazable. La salida forzada marca \texttt{requires\_human\_attention=true}, reforzando que el operador conserva la decisión final.

\section{Capa 3: explicación normativa con LLM, RAG y MCP}

La Capa 3 convierte la recomendación de Capa 2 en una salida útil para el operador. Produce \texttt{OperatorRecommendation}, que incluye prioridad, explicación textual, citas legales, resumen de reglas activadas, sugerencias no vinculantes y metadatos del LLM.

\subsection{LLM local}

El LLM local se utiliza para redactar explicaciones breves y consistentes. La inferencia se plantea con temperatura 0 para reducir variabilidad y mejorar reproducibilidad. El modelo no recibe el mandato de decidir prioridad; recibe la prioridad ya calculada, las reglas activadas y el contexto normativo recuperado.

Esta separación es clave para la seguridad del diseño. Un LLM puede redactar bien, pero no debe ser la autoridad decisoria en un dominio crítico. Si el LLM falla, la recomendación P1--P4 de Capa 2 sigue siendo válida y el sistema puede operar en modo degradado.

\subsection{RAG sobre corpus normativo CyL}

La recuperación aumentada conecta la explicación con documentos oficiales. El corpus incluye normativa estatal, autonómica, europea y fuentes de datos: Ley 17/2015, RD 524/2023, PLEGEM, Ley 4/2007 CyL, PLANCAL, INFOCAL, INUNCYL, MPCyL, Ley 11/2022, Registro 112 CyL, RGPD, LOPDGDD y Reglamento UE 2024/1689.

La función del RAG no es entrenar el modelo de prioridad. Su papel es proporcionar fragmentos normativos para que la explicación cite fundamentos relevantes. Esto reduce el riesgo de explicaciones genéricas y refuerza la trazabilidad del sistema.

\subsection{Herramientas MCP}

La Capa 3 expone tres herramientas MCP para la búsqueda normativa, la consulta del detalle de reglas y la construcción controlada de citas. La consulta meteorológica contextual queda diferida como ampliación futura. Estas funciones definen el catálogo básico de operaciones que el asistente puede realizar de forma controlada.

Estas herramientas permiten consultar el corpus, recuperar el detalle de reglas y construir citas legales. El uso de MCP separa la capacidad del LLM de las funciones verificables del sistema: buscar, citar y explicar reglas son operaciones con entradas y salidas controladas. Esta separación facilita la auditoría del proceso de generación de explicaciones a nivel de llamada API.

\subsection{Cache de sesión en Capa 3}

La Capa 3 incorpora una cache de sesión para evitar llamadas redundantes al LLM cuando se procesa el mismo texto operativo con la misma prioridad recomendada. La clave de cache se genera mediante SHA-256 a partir de la concatenación del texto del incidente y la prioridad: \texttt{sha256(texto || prioridad)}. Esta decisión reduce latencia en demostraciones y pruebas repetidas, sin alterar la recomendación de Capa 2 ni los contratos de salida.

El resultado cacheado es un objeto Pydantic \texttt{OperatorRecommendation}. Como ese objeto conserva el \texttt{incident\_id} del primer incidente que generó la explicación, la implementación clona el objeto mediante \texttt{model\_copy(update=\{...\})} cuando el \texttt{incident\_id} de la nueva inferencia es distinto. Esta copia evita colisiones en los logs de auditoría cruzada y mantiene la consistencia contractual: cada \texttt{InferenceLog} conserva hijos con el mismo identificador de incidente.

\section{Orquestador, fallback y latencia}

El backend actúa como orquestador. Recibe la entrada, valida el contrato, ejecuta Capa 1, invoca Capa 2, llama a Capa 3 si procede, registra latencias y devuelve la recomendación al cliente. También gestiona errores, timeouts y modo degradado.

\begin{figure}[htbp]
\centering
\small
\fbox{%
\begin{tabular}{c}
\texttt{IncidentInput} \\
\(\downarrow\) \\
\texttt{IncidentFeatures (Capa 1)} \\
\(\downarrow\) \\
\texttt{PriorityRecommendation = RuleFit-lite + Safety Gate (Capa 2)} \\
\(\downarrow\) \\
\texttt{OperatorRecommendation = RAG + LLM local + MCP / modo degradado (Capa 3)} \\
\(\downarrow\) \\
\texttt{OperatorDecision + InferenceLog}
\end{tabular}
}
\caption{Secuencia de inferencia desde la entrada del incidente hasta la decisión humana.}
\label{fig:secuencia-inferencia-cap6}
\end{figure}

El diseño contempla varios mecanismos de contingencia para mantener una respuesta del prototipo cuando se detectan fallos en componentes individuales. En concreto, el orquestador implementa las siguientes salvaguardas:

\begin{itemize}
    \item Si Capa 1 produce baja confianza, Capa 2 puede apoyarse en la línea base experta o en una respuesta conservadora.
    \item Si RuleFit no está disponible, se puede usar la línea base experta.
    \item Si Capa 3 supera el umbral temporal configurado o el LLM no está disponible, se devuelve una explicación estática basada en reglas activadas.
    \item Si la entrada contiene campos prohibidos, el sistema rechaza la solicitud en lugar de degradar silenciosamente.
\end{itemize}

Los objetivos internos de latencia se fijaron en un p95 igual o inferior a 500 ms para Capa 1 y Capa 2, y por debajo de 2 s p95 para el flujo completo con Capa 3. Son referencias de diseño para orientar la optimización, no un acuerdo de nivel de servicio del 112 ni garantías de despliegue productivo.

La latencia media observada en el entorno experimental con el LLM local es de \(4{,}898\) s por inferencia completa. Esta cifra supera el objetivo interno y muestra la limitación del hardware no dedicado. Como contingencia, si el modelo local no responde o supera el tiempo esperado, Capa 3 genera una explicación determinista basada en reglas activadas, probabilidades y plantillas, con latencia submilisegundo en las pruebas controladas. Así, la recomendación de Capa 2 no queda bloqueada por el componente generativo; este resultado acredita funcionamiento técnico, no cumplimiento de un acuerdo de nivel de servicio operativo.

\section{Trazabilidad y auditoría}

Cada inferencia genera un \texttt{InferenceLog}. El log incluye hash de entrada, snapshots de las salidas por capa, versiones de modelo, latencias, herramientas invocadas, timestamp y decisión del operador si existe. Esta información permite reconstruir por qué el sistema recomendó una prioridad y qué hizo finalmente la persona responsable. La persistencia de \texttt{OperatorDecision} y la conservación de snapshots por capa materializan los principios de supervisión humana y trazabilidad exigibles en sistemas de alto impacto, en coherencia con el Reglamento Europeo de Inteligencia Artificial \cite{EUAIAct2024}.

La trazabilidad sirve para tres fines. Primero, auditoría técnica: comprobar que el sistema usó los modelos y contratos esperados. Segundo, evaluación académica: medir errores, falsos negativos P1, calibración y divergencias operador-sistema. Tercero, alineación responsable del diseño: documentar supervisión humana, registro de decisiones y control de cambios, sin atribuir al prototipo una conformidad regulatoria certificada.

\section{Ejemplo de funcionamiento}

Considérese el incidente:

\begin{quote}
Varón inconsciente tras choque frontal en N-122, herido grave atrapado en el habitáculo. Han llamado varios testigos.
\end{quote}

La Capa 1 detecta atrapamiento, lesión grave, riesgo vital probable, avisos simultáneos y un accidente de tráfico. La Capa 2 recomienda P1 cuando RuleFit activa las reglas asociadas a atrapamiento y lesión grave. La Capa 3 recupera normativa relacionada con la protección de personas y PLANCAL, y redacta una explicación breve:

\begin{quote}
Prioridad P1 por indicios de riesgo vital. La descripción menciona herido grave y persona atrapada tras choque frontal, con varios avisos de testigos. La recomendación se basa en reglas de afectación a personas y requiere supervisión inmediata del operador.
\end{quote}

El operador puede aceptar la prioridad, modificarla o rechazarla. En cualquier caso, la divergencia queda registrada para evaluación posterior. Este registro histórico constituye una valiosa fuente de datos para el reentrenamiento y ajuste del modelo.

\section{Evaluación del diseño}

El Capítulo 9 mide si el diseño cumple sus objetivos. Las métricas principales de Capa 2 son F1 macro, sensibilidad P1, parsimonia del modelo, latencia, rendimiento temporal, análisis por subgrupos y matriz de confusión. El objetivo crítico es una sensibilidad P1 igual o superior a 0,85, porque los falsos negativos graves son el riesgo más importante. La calibración ECE queda identificada como mejora futura, y RuleFit debe respetar el límite de 30 reglas activas.

La Capa 3 se evaluará por validez contractual, latencia, presencia de citas en P1/P2 y fidelidad de explicación respecto a reglas activadas. La validación interna entre autores servirá para comparar etiquetas humanas, etiqueta débil y recomendación del sistema sobre un subconjunto de casos. Ambas líneas de evaluación permiten contrastar la coherencia de la arquitectura antes de planificar un piloto controlado.

\section{Conclusiones}

La arquitectura definida en este capítulo materializa la idea central del TFM: un sistema de apoyo a la decisión para incidentes 112 en Castilla y León que combina datos reales, normativa, aprendizaje interpretable, explicación local y supervisión humana. Esta propuesta busca superar las limitaciones de rigidez de los sistemas basados puramente en heurísticas. Asimismo, el diseño modular adoptado establece las bases técnicas para incorporar variables contextuales complejas en futuras versiones.

El diseño evita dos extremos poco defendibles. Por un lado, no se queda en reglas manuales rígidas: incorpora supervisión débil, RuleFit y evaluación cuantitativa. Por otro lado, no entrega la decisión a una caja negra ni a un LLM generativo: mantiene la prioridad en un núcleo interpretable y reserva el LLM para explicar y citar. Esta posición intermedia es la que permite aspirar a un TFM excelente: una propuesta ambiciosa, pero acotada; técnica, pero auditable; innovadora, pero prudente en un dominio crítico.

El siguiente capítulo desarrolla el diseño de datos, la auditoría del Registro 112 CyL, la construcción de etiquetas académicas y las garantías anti-\textit{leakage} que alimentan la arquitectura aquí descrita. Este análisis empírico es fundamental para comprender cómo se entrenan y validan los componentes predictivos de la Capa 2. Con ello se asegura el rigor científico y la reproducibilidad de todo el proceso experimental del prototipo.
